\chapterstart{Качество обучения и проверка повторяемости}
\label{sec:training}
\subsection{Какие решения найдены в основной серии}
Все 15 обучений основной серии завершены, каждое содержит 100\,000 переходов и 99\,872 обновления. В Таблице~\ref{tab:main} режим оценки совпадает с режимом обучения, а значения относятся только к двадцати валидационным состояниям; SAC-off, в частности, не оценивается здесь с незаявленной дополнительной защитой.
\begin{table}[htbp]\caption{Основная серия: успехи лучших / последних весов из 20}\label{tab:main}
\begin{smalltable}\begin{tabularx}{\linewidth}{@{}X*{3}{>{\centering\arraybackslash}p{26mm}}@{}}\toprule
Конфигурация & Seed 0 & Seed 1 & Seed 2\\\midrule
SAC-on & 20 / 20 & 20 / 20 & 20 / 20\\
TQC-on & 20 / 20 & 20 / 20 & 20 / 20\\
DQN-on & 20 / 11 & 14 / 14 & 20 / 20\\
DDPG-on & 20 / 20 & 0 / 0 & 0 / 0\\
SAC-off & 0 / 0 & 0 / 0 & 0 / 0\\\midrule
Классика & \multicolumn{3}{c}{20/20 при off и on; без обучения}\\\bottomrule
\end{tabularx}\end{smalltable}
\source{Результаты вычислительных экспериментов автора; валидационные состояния, именованные исключены.}
\end{table}

SAC-on и TQC-on находят успешные лучшие модели и сохраняют успех последних весов во всех трёх запусках. Метрика при этом насыщена — 20 из 20 у обоих методов и всех seeds, — поэтому безусловное превосходство одного над другим по этим данным установить нельзя. DQN и особенно DDPG заметнее зависят от seed: у DQN seed 0 последние веса дают 11 успехов против 20 у выбранных. Исходный SAC-off не решает задачу ни в одном запуске, и это результат конкретного бюджета и настроек, а не невозможность управления без фильтра.

Три именованных сценария рассматриваются отдельно и в отбор не входят. Для best SAC-on и TQC-on каждый seed даёт 3/3; DQN-on — 3/3, 2/3, 3/3; DDPG-on — 3/3, 0/3, 0/3; SAC-off — нули во всех трёх. Вместе с двадцатью валидационными состояниями это даёт описательные суммы 23/23 и 16/23, которые иногда удобнее для краткости, однако знаменателем выбора остаются именно 20 состояний.

Кривые на Рисунке~\ref{fig:learning} показывают, что одинаковый бюджет приводит к весьма разным траекториям качества. Каждая точка — оценка сохранённой модели; отрезок между точками лишь соединяет соседние измерения, и о качестве между ними ничего не известно. Потеря качества к концу обучения, хорошо заметная у DQN seed 0, и объясняет, зачем лучшие и последние веса хранятся раздельно.
\fig{main_learning}{Успех на двадцати валидационных состояниях в основной серии. Фильтр при периодической оценке соответствует режиму обучения; каждый seed показан отдельно.}{.89}{fig:learning}

\subsection{Два разных сравнения фильтра}
Фильтр участвует в эксперименте дважды — при обучении и при исполнении, — и эти два его вклада нужно различать. Строки Таблицы~\ref{tab:main} сравнивают \emph{процедуры целиком}: SAC-on и SAC-off имеют одинаковые базовые настройки и одинаковые seeds, но различаются и режимом обучения, и режимом оценки. Разрыв 20/20 против 0/20 на каждом seed относится именно к паре целых процедур.

Чтобы отделить режим исполнения, достаточно уравнять его. Сохранённые веса SAC-off, исполненные \emph{с} фильтром, дают те же 0/20 на всех трёх seeds, то есть выдача защиты при оценке эту процедуру не спасает. В обратную сторону у SAC-on исполнение без фильтра снижает результат seed 1 с 20/20 до 0/20, тогда как seeds 0 и 2 сохраняют 20/20. Режим исполнения, следовательно, влияет на исход, но наблюдаемый разрыв между процедурами им не объясняется.

Остаётся вклад фильтра при обучении, и здесь данные позволяют меньше. Наличие защиты меняет не только правило исполнения, но и весь собранный опыт: длительности эпизодов, посещаемые состояния и дальнейшую последовательность случайных чисел. Эти механизмы в основной серии не разделялись, поэтому корректное утверждение звучит так: процедура обучения с фильтром в данном протоколе превосходит сопоставимую процедуру без него, а какая доля эффекта приходится на саму защиту, а какая на изменившийся опыт, из этих данных не следует. Универсальным свойством SAC результат тоже не является.

Второй способ исполняет \emph{одни и те же веса} с фильтром и без него. У best SAC-on seed 1, TQC-on seeds 0 и 1, DDPG-on seed 0 и всех трёх DQN-on выключение защиты снижает validation-успех до нуля. У best SAC-on seeds 0 и 2, а также TQC-on seed 2 фильтр на этих стартах вообще не вмешивается, и оба режима дают 20/20. Полная матрица приведена в Приложении~\ref{app:main}. Зависимость исполнения от защитного слоя отсюда видна прямо; при этом даже нулевая доля вмешательств при оценке не означает, что фильтр не повлиял на обучение, которое эти веса породило.

У классического контроллера соотношение иное. Оба режима успешны в 23/23 описательных сценариях, но при off в 21 из них превышаются желаемые 0,24 м, хотя события на 0,25 м не происходит. Один и тот же исходный успех, таким образом, допускает разные характеристики безопасности. На 713 внешних on-слотах основной серии аудит не нашёл ни превышений 0,24 м, ни событий положения, ни отказов фильтра — но эти слоты включают разные метки best/last и повторные старты, поэтому 713 независимых испытаний надёжности из них не получится.

\subsection{Ограниченный подбор: победа по правилу и границы объяснения}
Stage 1 содержит 12 новых обучений: DDPG-on и SAC-off, по два однофакторных кандидата и три значения seed. Таблица~\ref{tab:stage} даёт результаты best в режиме соответствующего обучения; базовый вариант взят из основной серии без повторного обучения.
\begin{table}[htbp]\caption{Stage 1: успехи best из 20, seeds 0 / 1 / 2}\label{tab:stage}
\begin{smalltable}\begin{tabularx}{\linewidth}{@{}Xcc@{}}\toprule
Настройка & DDPG-on, оценка on & SAC-off, оценка off\\\midrule
Базовая настройка & 20 / 0 / 0 & 0 / 0 / 0\\
Скорость обучения $10^{-4}$ & 0 / 13 / 0 & 0 / 0 / 0\\
Начальный сбор 5000 переходов & 0 / 20 / 20 & 20 / 0 / 0\\\bottomrule
\end{tabularx}\end{smalltable}
\source{Результаты вычислительных экспериментов автора; валидационные состояния, именованные исключены.}
\end{table}

По правилу~\eqref{eq:config} строгим победителем в обеих группах оказывается вариант с начальным сбором 5000 переходов (warmup5000). Существенно, где именно достигается этот выигрыш: минимальный успех по seeds остаётся нулевым, и решает вторая компонента ключа — средняя доля, которая у DDPG растёт с $1/3$ до $2/3$, а у SAC-off с 0 до $1/3$. Вариант $10^{-4}$ базовую настройку не улучшает. Точные ключи сохранены в электронных материалах, округлённые приведены в Приложении~\ref{app:selection}. Называть после этого DDPG или SAC-off устойчивыми к seed оснований нет.

За средними значениями скрывается заметный разброс. У DDPG warmup5000 seed 1 успех достигает 20/20 на 80 и 90 тыс. переходов, а затем падает до 0/20 на 100 тыс.; seed 2 после промежуточных провалов удерживает 20/20 на 80–100 тыс. У SAC-off весь выигрыш сосредоточен в seed 0, где best даёт 20/20, а last — 17/20 при off, тогда как seeds 1 и 2 двухсекундного удержания не приобретают вовсе. Одна улучшившаяся средняя величина, как видно, всю процедуру не описывает.

Гипотеза, стоявшая за увеличенным начальным сбором, предполагала улучшение раннего исследования верхней области. Проверка дала расщеплённый ответ: к 20 тыс. переходов доля таких образцов улучшилась у всех трёх запусков DDPG, но только у одного из трёх запусков SAC-off, тогда как заранее требовалось не менее двух. Следовательно, \emph{победа SAC-off warmup5000 по итоговому ключу не подтверждает заявленный ранний механизм}. Для DDPG признаки гипотезу поддерживают, но не изолируют её: вместе с warmup меняются и данные, и случайная последовательность, и число обновлений. Переоценивание критика, недостаточная сеть или нехватка переходов как причины остаются недоказанными.

\subsection{Подтверждение на новых значениях seed}
Все шесть новых обучений warmup5000 и шесть внешних оценок завершены; каждое обучение содержит 100\,000 переходов и 95\,000 обновлений. При неизменных параметрах best DDPG-on на seeds 3/4/5 даёт 4/0/20 успехов из 20, а last — 0/0/0. SAC-off даёт нули и у лучших, и у последних весов во всех трёх случаях. Ни одна из конфигураций не проходит закреплённый критерий: не каждый best достигает 16/20, и нет двух успешных last.

Рисунок~\ref{fig:confirmation} сохраняет различие между seeds отбора и seeds проверки. Отрицательный результат получен в корректно завершённой серии и пропущенными заданиями не вызван. Правильность применения правила Stage 1 он не отменяет — он ограничивает практическую ценность выбранного улучшения. Дополнительный бюджет по итогам этой проверки автоматически не назначался.
\fig{confirmation}{Warmup5000 на seeds 0–5: лучшие и последние веса на двадцати валидационных состояниях в режиме обучения. Вертикальная линия отделяет Stage 1 от confirmation; горизонтальный пунктир — порог 16/20.}{.96}{fig:confirmation}

\subsection{Маятник сверху: почему это ещё не успех}
Расхождение награды и успеха прямо зафиксировано в журналах last DDPG-on confirmation для seeds 3 и 5. Там все двадцать validation-эпизодов в последние две секунды удовлетворяют условиям на угол, угловую скорость и координату из~\eqref{eq:hold} — маятник действительно стоит вверху. Прерывает удержание единственная оставшаяся компонента: скорость тележки выходит за 0,20 м/с. Из 4020 отсчётов последних двух секунд скорость допустима в 1925 (47,89\%) для seed 3 и в 2404 (59,80\%) для seed 5. Это доли отсчётов по двадцати траекториям, не доли успешных эпизодов; нарушение компоненты критерия зафиксировано здесь по журналам, а не по визуальному впечатлению от анимации.

Расхождение награды и успеха здесь тоже измеримо. Для seed 3 средний возврат от best к last растёт с 273,68 до 706,05, тогда как успех падает с 4/20 до 0/20, а среднее заключительное удержание — с 0,8895 до 0,3325 с. У seed 5 успех падает с 20/20 до 0/20, а удержание — с 7,8915 до 0,0115 с. Конкретная пара сохранённых траекторий seed 5 показана на Рисунке~\ref{fig:ddpg}; в этом отдельном примере возврат, наоборот, снижается, и смешивать два наблюдения не следует. Синхронное воспроизведение той же пары доступно в \vref{02_ddpg_best_vs_last}{видеоприложении 2}: последние веса визуально держат маятник около верха, но допустимой скорости тележки не обеспечивают.
\fig{ddpg_criterion}{DDPG-on warmup5000, seed 5, validation\_1000, оценка on: best 90k и last 100k. На заключительных двух секундах угол допустим, но last пересекает порог скорости. График построен по сохранённым состояниям двух эпизодов.}{.94}{fig:ddpg}

Такое расхождение прямо следует из определений. При $v=0{,}30$ м/с множитель $q(v/2)\approx0{,}978$ ещё почти не отличается от единицы, хотя жёсткий порог скорости удержания уже нарушен; к тому же возврат суммируется по всей траектории, а успех зависит только от её конца. Различие метрик этим объясняется, но причина изменения весов актора или критика — нет. У других неудачных запусков и траекторий механизмы могут быть иными, и отдельный ролик общего механизма ухудшения DDPG не доказывает.

Результат главы двойственен. В основной серии получены работающие SAC и TQC с фильтром, однако качество обучения воспроизводится не у всех методов одинаково, а настройка, выигравшая отбор на этапе разработки, проверку на новых запусках не прошла. Всё это говорит об изменчивости \emph{обучения}. Следующая глава задаёт другой вопрос — насколько чувствительны уже обученные и закреплённые политики к начальному состоянию.
